\documentclass[aspectratio=169]{beamer}

% Tema UFS --- Universidade Federal de Sergipe
\usetheme{UFS}

% Pacotes base
\usepackage[utf8]{inputenc}
\usepackage[portuguese]{babel}
\usepackage{amsmath,amssymb}
\usepackage{graphicx}
\usepackage{booktabs}
\usepackage{xcolor}

% TikZ e bibliotecas
\usepackage{tikz}
\usetikzlibrary{shapes.geometric, shapes.multipart, arrows.meta, positioning,
  matrix, fit, backgrounds, decorations.pathreplacing, patterns, calc}

% Listagens --- paleta VS Code Light+
\usepackage{listings}
\definecolor{bgcolor}{HTML}{FFFFFF}
\definecolor{linecolor}{HTML}{DDDDDD}
\definecolor{numcolor}{HTML}{999999}
\definecolor{kwcolor}{HTML}{0000FF}
\definecolor{strcolor}{HTML}{A31515}
\definecolor{cmtcolor}{HTML}{008000}

\lstdefinestyle{ufsStyle}{
  backgroundcolor=\color{bgcolor},
  basicstyle=\ttfamily\footnotesize\color{black},
  keywordstyle=\color{kwcolor}\bfseries,
  stringstyle=\color{strcolor},
  commentstyle=\color{cmtcolor}\itshape,
  numberstyle=\tiny\color{numcolor},
  numbers=left, numbersep=12pt,
  xleftmargin=20pt, framexleftmargin=20pt,
  breaklines=true, tabsize=4,
  keepspaces=true, showspaces=false,
  showstringspaces=false, showtabs=false,
  frame=single, rulecolor=\color{linecolor},
  framesep=4pt, captionpos=b, upquote=true,
  columns=fullflexible,
}
\lstset{style=ufsStyle, language=C}

\lstdefinestyle{tinyC}{
  style=ufsStyle, language=C, numbers=none,
  basicstyle=\ttfamily\scriptsize\color{black},
}
\lstdefinestyle{codC}{
  style=tinyC, basicstyle=\ttfamily\fontsize{7.5}{8.6}\selectfont\color{black},
}
\lstdefinestyle{shellStyle}{
  style=ufsStyle, language=bash, numbers=none,
  basicstyle=\ttfamily\scriptsize\color{black},
}

% --- Estilos TikZ reutilizados ---
\tikzset{
  no/.style={rectangle split, rectangle split horizontal, rectangle split parts=2,
    draw=UFSBlue, fill=UFSBlue!8, minimum height=0.55cm,
    font=\scriptsize\ttfamily, inner sep=3pt,
    rectangle split part fill={UFSBlue!8, white}},
  nonovo/.style={no, draw=UFSRed, rectangle split part fill={UFSRed!8, white}},
  nolivre/.style={no, draw=UFSGray, dashed, rectangle split part fill={black!5, black!5}},
  var/.style={rectangle, draw=UFSGray, fill=black!5, minimum width=0.7cm,
    minimum height=0.55cm, font=\scriptsize\ttfamily},
  varr/.style={var, draw=UFSRed, fill=UFSRed!8},
  nulo/.style={font=\scriptsize\ttfamily, text=UFSGray},
  rot/.style={font=\scriptsize\ttfamily, text=UFSGray},
  leg/.style={font=\scriptsize, text=UFSGray},
  seta/.style={-{Stealth[length=2mm]}, thick, UFSBlue},
  setar/.style={-{Stealth[length=2mm]}, thick, UFSRed},
  setag/.style={-{Stealth[length=2mm]}, thick, UFSGray},
}

% Centro da parte "prox" de um no (ponto de partida da seta)
\newcommand{\ptr}[1]{($(#1.one split north)!0.5!(#1.south east)$)}

% Referencia da fonte no rodape do frame
\newcommand{\fonte}[1]{\vfill\hfill{\scriptsize\textcolor{UFSGray}{Fonte: #1}}}

% --- Metadados ---
\title{Lista encadeada: busca, inserção e remoção}
\subtitle{Achar, pôr e tirar um nó em qualquer ponto da lista}
\author{Prof.\ André Yoshiaki Kashiwabara}
\institute{DCOMP -- Universidade Federal de Sergipe\\
\texttt{andreyoshiaki@dcomp.ufs.br}}
\date{}

\begin{document}

\begin{frame}[plain]
  \titlepage
\end{frame}

\begin{frame}{O que você vai aprender hoje}
  \begin{center}
  \begin{tikzpicture}[
    item/.style={rectangle, draw=UFSBlue, fill=UFSBlue!10, rounded corners,
      minimum width=10.4cm, minimum height=0.55cm, font=\small}
  ]
    \node[item] (t1) at (0, 0)    {Buscar: andar até achar --- ou até o \texttt{NULL}};
    \node[item] (t2) at (0,-0.72) {Inserir em qualquer posição: a lista em ordem};
    \node[item] (t3) at (0,-1.44) {Remover: desligar o nó sem quebrar a corrente};
    \node[item] (t4) at (0,-2.16) {Ponteiro para ponteiro: um código sem casos especiais};
    \node[item] (t5) at (0,-2.88) {Quanto custa cada operação};
    \node[item] (t6) at (0,-3.60) {Exercícios};
    \draw[seta] (t1) -- (t2);
    \draw[seta] (t2) -- (t3);
    \draw[seta] (t3) -- (t4);
    \draw[seta] (t4) -- (t5);
    \draw[seta] (t5) -- (t6);
  \end{tikzpicture}
  \end{center}
\end{frame}

% ==================================================================
\section{A pergunta de hoje}
% ==================================================================

\begin{frame}{O problema: uma lista que não é só uma pilha de valores}
  \begin{columns}[T]
    \begin{column}{0.5\textwidth}
      \textbf{Até aqui}\\[0.3em]
      \small
      A \texttt{lista.c} da aula passada só sabia inserir \emph{no início},
      percorrer e liberar. A ordem saía invertida, e não havia como
      perguntar ``o 20 está aí?'' nem tirar um valor do meio.
    \end{column}
    \begin{column}{0.5\textwidth}
      \textbf{O que vamos descobrir}\\[0.3em]
      \small
      Para achar, pôr e tirar em qualquer ponto, basta andar pelas setas
      --- mas cada operação esconde um \textbf{caso especial} que derruba
      o programa se for esquecido.
    \end{column}
  \end{columns}
  \begin{block}{A pergunta de hoje}
    Como buscar, inserir e remover em qualquer posição de uma lista
    encadeada --- e quais casos especiais aparecem no caminho?
  \end{block}
\end{frame}

\begin{frame}[fragile]{O programa de hoje: \texttt{lista.c}, versão 2 (1/3)}
  \begin{columns}[T]
    \begin{column}{0.44\textwidth}
\begin{lstlisting}[style=codC]
/* lista.c --- lista em ordem crescente:
   busca, insercao e remocao */
#include <stdio.h>
#include <stdlib.h>

typedef struct no {
    int valor;
    struct no *prox;
} No;
\end{lstlisting}
      \small \texttt{imprime} e \texttt{libera} continuam iguais aos da aula
      passada.
    \end{column}
    \begin{column}{0.52\textwidth}
\begin{lstlisting}[style=codC]
/* devolve o primeiro no com valor x,
   ou NULL */
No *busca(No *ini, int x)
{
    No *p = ini;
    while (p != NULL && p->valor != x)
        p = p->prox;
    return p;
}
\end{lstlisting}
    \end{column}
  \end{columns}
\end{frame}

\begin{frame}[fragile]{\texttt{lista.c} (2/3): inserir em ordem e remover}
  \vspace{-0.5em}
  \begin{columns}[T]
    \begin{column}{0.5\textwidth}
\begin{lstlisting}[style=codC]
No *insere_ordenado(No *ini, int x)
{
    No *novo = malloc(sizeof *novo);
    if (novo == NULL) { perror("malloc"); exit(1); }
    novo->valor = x;
    if (ini == NULL || x < ini->valor) {
        novo->prox = ini;  /* 1o */
        return novo;
    }
    No *ant = ini;
    while (ant->prox != NULL &&
           ant->prox->valor < x)
        ant = ant->prox;
    novo->prox = ant->prox;
    ant->prox = novo;
    return ini;
}
\end{lstlisting}
    \end{column}
    \begin{column}{0.46\textwidth}
\begin{lstlisting}[style=codC]
No *remove_valor(No *ini, int x)
{
    No *ant = NULL, *p = ini;
    while (p != NULL && p->valor != x) {
        ant = p;
        p = p->prox;
    }
    if (p == NULL)     /* nao achou */
        return ini;
    if (ant == NULL)   /* 1o */
        ini = p->prox;
    else
        ant->prox = p->prox;
    free(p);
    return ini;
}
\end{lstlisting}
    \end{column}
  \end{columns}
\end{frame}

\begin{frame}[fragile]{\texttt{lista.c} (3/3): o \texttt{main}}
  \vspace{-0.5em}
  \begin{columns}[c]
    \begin{column}{0.62\textwidth}
\begin{lstlisting}[style=codC]
int main(void)
{
    int entrada[] = {30, 10, 40, 20};
    No *lista = NULL;
    for (int i = 0; i < 4; i++)
        lista = insere_ordenado(lista, entrada[i]);
    imprime(lista);
    printf("busca 20: %s\n", busca(lista, 20) ? "achou" : "nao achou");
    printf("busca 25: %s\n", busca(lista, 25) ? "achou" : "nao achou");
    lista = remove_valor(lista, 10);
    lista = remove_valor(lista, 40);
    lista = remove_valor(lista, 99);
    imprime(lista);
    libera(lista);
    return 0;
}
\end{lstlisting}
    \end{column}
    \begin{column}{0.34\textwidth}
      \small
      Os valores chegam fora de ordem: 30, 10, 40, 20.

      \medskip
      Depois de duas buscas, o programa tenta remover três valores --- um
      deles não está na lista.
    \end{column}
  \end{columns}
\end{frame}

\begin{frame}{O que este programa imprime?}
  \begin{block}{Antes de compilar, escreva sua resposta}
    \begin{enumerate}
      \item Em que ordem os valores aparecem na primeira linha? Compare com a
        aula passada.
      \item O que as duas buscas respondem?
      \item Qual é a última linha impressa?
      \item Dos três \texttt{remove\_valor}, qual passa pelo
        \texttt{if (ant == NULL)}? E qual não chega a chamar
        \texttt{free}?
    \end{enumerate}
  \end{block}
\end{frame}

\begin{frame}[fragile]{Compile, execute e compare}
\begin{lstlisting}[style=shellStyle]
$ gcc -Wall -Wextra lista.c -o lista
$ ./lista
10 -> 20 -> 30 -> 40 -> NULL
busca 20: achou
busca 25: nao achou
20 -> 30 -> NULL
\end{lstlisting}
  \begin{columns}[T]
    \begin{column}{0.48\textwidth}
      \begin{block}{O que mudou}
        \small A lista saiu \textbf{em ordem}, embora os valores tenham
        chegado embaralhados: cada um foi encaixado no lugar certo.
      \end{block}
    \end{column}
    \begin{column}{0.48\textwidth}
      \small Remover o 10 mudou o início; remover o 99 não fez nada. Cada
      função trata esses casos com um \texttt{if} --- e é deles que esta
      aula trata.
    \end{column}
  \end{columns}
\end{frame}

% ==================================================================
\section{Buscar}
% ==================================================================

\begin{frame}{Definição: busca sequencial}
  \begin{block}{Busca sequencial numa lista}
    \small Começando no início, compare cada nó com o valor procurado e siga
    o \texttt{prox}. Pare ao \textbf{achar} (devolva o nó) ou ao chegar ao
    \textbf{\texttt{NULL}} (devolva \texttt{NULL}: não está na lista).
  \end{block}
  \begin{center}
  \begin{tikzpicture}[node distance=0.85cm]
    \node[var] (ini) {lista};
    \node[no, right=0.7cm of ini] (a) {10 \nodepart{two} \phantom{xx}};
    \node[no, right=of a] (b) {20 \nodepart{two} \phantom{xx}};
    \node[nonovo, right=of b] (c) {30 \nodepart{two} \phantom{xx}};
    \node[no, right=of c] (d) {40 \nodepart{two} \phantom{xx}};
    \node[nulo, right=0.6cm of d] (n) {NULL};
    \draw[seta] (ini) -- (a);
    \draw[seta] \ptr{a} -- (b);
    \draw[seta] \ptr{b} -- (c);
    \draw[seta] \ptr{c} -- (d);
    \draw[seta] \ptr{d} -- (n);
    \foreach \q/\t in {a/1, b/2, c/3} {
      \node[rot, above=0.12cm of \q] {p (\t)};
    }
    \node[leg, text=UFSRed, right=0.3cm of n] {busca 30: 3 passos};
  \end{tikzpicture}
  \end{center}
  \begin{exampleblock}{Intuição}
    \small Não há índice: a lista só se deixa ler na ordem das setas. Para
    saber que 25 \emph{não} está lá, é preciso olhar todos os nós.
  \end{exampleblock}
  \fonte{Cormen et al.\ (2009), §10.2, \textsc{List-Search}~\cite{cormen2009introduction};
    Celes et al.\ (2004), cap.~10~\cite{celes2004introducao}}
\end{frame}

\begin{frame}[fragile]{E se a condição vier na outra ordem?}
  \begin{columns}[c]
    \begin{column}{0.52\textwidth}
\begin{lstlisting}[style=tinyC]
/* busca_errada.c */
No *busca(No *ini, int x)
{
    No *p = ini;
    while (p->valor != x && p != NULL)
        p = p->prox;
    return p;
}
\end{lstlisting}
      \small A mesma condição, só que com os dois testes trocados de lugar.
      O \texttt{\&\&} não é comutativo?
    \end{column}
    \begin{column}{0.44\textwidth}
      \begin{block}{Pense antes de passar o slide}
        \small
        \begin{enumerate}
          \item \texttt{busca(lista, 20)} ainda funciona?
          \item E \texttt{busca(lista, 25)}? Siga \texttt{p} até o fim.
          \item Faz diferença compilar com \texttt{-O2}?
        \end{enumerate}
      \end{block}
    \end{column}
  \end{columns}
\end{frame}

\begin{frame}[fragile]{Execute e compare}
  \vspace{-0.5em}
\begin{lstlisting}[style=shellStyle]
$ gcc -Wall -Wextra busca_errada.c -o busca_errada     # nenhum aviso
$ ./busca_errada
10 -> 20 -> 30 -> 40 -> NULL
busca 20: achou
Segmentation fault: 11
$ gcc -O2 -Wall -Wextra busca_errada.c -o busca_errada-O2
$ ./busca_errada-O2
10 -> 20 -> 30 -> 40 -> NULL
busca 20: achou
busca 25: achou
20 -> 30 -> NULL
\end{lstlisting}
  \vspace{-0.6em}
  \centering\scriptsize\textcolor{UFSGray}{Sem otimização, o programa cai. Com
  \texttt{-O2}, ele não cai --- e afirma que 25 está na lista.}
\end{frame}

\begin{frame}{A resposta: o \texttt{\&\&} para no primeiro falso}
  \begin{columns}[T]
    \begin{column}{0.48\textwidth}
      \small
      \texttt{a \&\& b} avalia \texttt{a} primeiro; se \texttt{a} é falso,
      \texttt{b} \textbf{nem é avaliado} (curto-circuito). A ordem dos
      testes é parte do significado.

      \medskip
      Na busca por 25, \texttt{p} chega a \texttt{NULL}. Na versão certa,
      \texttt{p != NULL} falha e \texttt{p->valor} nunca é lido. Na
      trocada, \texttt{p->valor} lê o endereço 0: \emph{segmentation
      fault}.
    \end{column}
    \begin{column}{0.48\textwidth}
      \begin{alertblock}{E o \texttt{-O2}?}
        \small Ler \texttt{p->valor} com \texttt{p == NULL} é comportamento
        indefinido. O compilador pode supor que isso nunca acontece --- e,
        se \texttt{p} nunca é \texttt{NULL}, o teste \texttt{p != NULL} e o
        ``não achou'' podem ser descartados.
      \end{alertblock}
      \small Proteja sempre o acesso: \texttt{p != NULL \&\&
      p->\ldots} --- nessa ordem.
    \end{column}
  \end{columns}
  \fonte{ISO/IEC 9899:2018, §6.5.13 e §3.4.3~\cite{iso2018c17};
    Kernighan \& Ritchie (1989), §2.6~\cite{kernighan1989c};
    medido com \texttt{busca\_errada.c}}
\end{frame}

\begin{frame}{Recapitulando: buscar}
  \begin{block}{O que aprendemos}
    \begin{itemize}
      \item buscar é percorrer até achar o valor ou chegar ao \texttt{NULL};
      \item a função devolve o nó encontrado --- ou \texttt{NULL};
      \item \texttt{p != NULL} vem \emph{antes} de \texttt{p->valor}: o
        \texttt{\&\&} para no primeiro falso;
      \item com um ponteiro nulo, o programa pode cair --- ou, otimizado,
        dar uma resposta errada sem aviso.
    \end{itemize}
  \end{block}
\end{frame}

% ==================================================================
\section{Inserir em qualquer posição}
% ==================================================================

\begin{frame}[fragile]{Onde encaixar o 25?}
  \begin{columns}[c]
    \begin{column}{0.5\textwidth}
\begin{lstlisting}[style=tinyC]
No *ant = ini;
while (ant->prox != NULL &&
       ant->prox->valor < x)
    ant = ant->prox;
novo->prox = ant->prox;
ant->prox = novo;
\end{lstlisting}
      \small A lista é \texttt{10 -> 20 -> 30 -> 40} e \texttt{x} vale 25.
    \end{column}
    \begin{column}{0.46\textwidth}
      \begin{block}{Pense antes de passar o slide}
        \small
        \begin{enumerate}
          \item Em que nó \texttt{ant} para?
          \item Por que o laço olha \texttt{ant->prox->valor}, um nó
            \emph{à frente}, e não o próprio \texttt{ant}?
          \item E se \texttt{x} for 50? E 5?
        \end{enumerate}
      \end{block}
    \end{column}
  \end{columns}
\end{frame}

\begin{frame}{A resposta: parar no nó \emph{anterior}}
  \begin{columns}[c]
    \begin{column}{0.58\textwidth}
  \begin{tikzpicture}[node distance=0.55cm]
    \node[no] (a) {10 \nodepart{two} \phantom{xx}};
    \node[no, right=of a] (b) {20 \nodepart{two} \phantom{xx}};
    \node[no, right=1.3cm of b] (c) {30 \nodepart{two} \phantom{xx}};
    \node[no, right=of c] (d) {40 \nodepart{two} \phantom{xx}};
    \node[nulo, right=0.4cm of d] (n) {NULL};
    \node[varr, above=0.45cm of b] (ant) {ant};
    \draw[setar] (ant) -- (b);
    \draw[seta] \ptr{a} -- (b);
    \draw[setag, dashed] \ptr{b} -- (c);
    \draw[seta] \ptr{c} -- (d);
    \draw[seta] \ptr{d} -- (n);
    \node[nonovo] (x) at ($(b.east)!0.5!(c.west)+(0,-1.25)$) {25 \nodepart{two} \phantom{xx}};
    \draw[setar] \ptr{x} to[out=0, in=-90] node[right, leg, text=UFSRed, pos=0.6] {1} (c.south);
    \draw[setar] \ptr{b} to[out=-60, in=180] node[left, leg, text=UFSRed, pos=0.4] {2} (x.west);
    \node[leg, anchor=west] at (-0.3,-2.2) {\textbf{1.} \texttt{novo->prox = ant->prox;}\quad \textbf{2.} \texttt{ant->prox = novo;}};
  \end{tikzpicture}
    \end{column}
    \begin{column}{0.38\textwidth}
      \small
      A lista só tem setas para a frente: ao chegar ao 30, já não dá para
      voltar ao 20 e mudar o \texttt{prox} dele. Por isso o laço espia um
      nó à frente e para no \textbf{anterior}.

      \medskip
      A ordem 1--2 é a mesma da aula passada: ligar o novo antes de
      desviar a seta antiga.
    \end{column}
  \end{columns}
  \fonte{Celes et al.\ (2004), cap.~10~\cite{celes2004introducao};
    Feofiloff (2009), cap.~4~\cite{feofiloff2009algoritmos}}
\end{frame}

\begin{frame}{Quatro situações, dois caminhos no código}
  \begin{center}
    \small
    \begin{tabular}{@{}llll@{}}
      \toprule
      situação & exemplo & quem trata & muda o início?\\
      \midrule
      lista vazia & 25 em \texttt{NULL} & \texttt{if (ini == NULL || \ldots)} & sim\\
      antes do primeiro & 5 em \texttt{10 -> \ldots} & \texttt{if (\ldots\ || x < ini->valor)} & sim\\
      no meio & 25 & laço + passos 1 e 2 & não\\
      no fim & 50 & laço + passos 1 e 2 & não\\
      \bottomrule
    \end{tabular}
  \end{center}
  \small
  No fim, o laço para no último nó porque \texttt{ant->prox == NULL}; o
  passo 1 copia esse \texttt{NULL} para o novo nó, que vira o último
  \emph{sem} caso especial.

  \medskip
  Só o primeiro nó precisa de tratamento à parte: \textbf{não existe nó
  anterior} a ele --- quem aponta para o primeiro é a variável
  \texttt{lista} do \texttt{main}.
  \fonte{Celes et al.\ (2004), cap.~10~\cite{celes2004introducao};
    King (2008), §17.5~\cite{king2008c}}
\end{frame}

\begin{frame}{E inserir sempre no fim, quanto custa?}
  \begin{columns}[c]
    \begin{column}{0.52\textwidth}
      \begin{center}
      \small
      \begin{tabular}{@{}rrr@{}}
        \toprule
        $n$ & no início & no fim\\
        \midrule
        10\,000  & 0,0003\,s & 0,058\,s\\
        20\,000  & 0,0002\,s & 0,147\,s\\
        40\,000  & 0,0004\,s & 0,625\,s\\
        80\,000  & 0,0007\,s & 2,479\,s\\
        160\,000 & 0,0029\,s & 10,27\,s\\
        \bottomrule
      \end{tabular}

      \medskip
      \scriptsize compilado com \texttt{-O2}
      \end{center}
    \end{column}
    \begin{column}{0.44\textwidth}
      \small
      Para pôr no fim, \texttt{insere\_fim} percorre a lista inteira a
      cada chamada: $0 + 1 + \cdots + (n-1) \approx n^2/2$ passos.

      \medskip
      Dobrou $n$, o tempo \textbf{quadruplicou} --- o mesmo $n^2$ que o
      vetor pagava para inserir no começo.

      \medskip
      Guardar também um ponteiro para o \textbf{último} nó elimina a
      caminhada (Exercício 1).
    \end{column}
  \end{columns}
  \fonte{Cormen et al.\ (2009), §10.2~\cite{cormen2009introduction};
    medido com \texttt{fim\_vs\_inicio.c}}
\end{frame}

\begin{frame}{Recapitulando: inserir em qualquer posição}
  \begin{block}{O que aprendemos}
    \begin{itemize}
      \item para inserir, é preciso parar no nó \emph{anterior} à posição;
      \item dois ponteiros, nesta ordem: \texttt{novo->prox = ant->prox;
        ant->prox = novo;}
      \item o fim não é especial; o início é, porque não tem anterior;
      \item inserir no fim percorrendo a lista custa $n$ passos por
        inserção.
    \end{itemize}
  \end{block}
\end{frame}

% ==================================================================
\section{Remover}
% ==================================================================

\begin{frame}{Remover: desviar a seta, depois liberar}
  \begin{columns}[c]
    \begin{column}{0.56\textwidth}
  \begin{tikzpicture}[node distance=0.65cm]
    \node[no] (a) {10 \nodepart{two} \phantom{xx}};
    \node[nolivre, right=of a] (b) {20 \nodepart{two} \phantom{xx}};
    \node[no, right=of b] (c) {30 \nodepart{two} \phantom{xx}};
    \node[nulo, right=0.45cm of c] (n) {NULL};
    \node[var, above=0.45cm of a] (ant) {ant};
    \node[varr, above=0.45cm of b] (p) {p};
    \draw[setag] (ant) -- (a);
    \draw[setar] (p) -- (b);
    \draw[setag, dashed] \ptr{a} -- (b);
    \draw[setag, dashed] \ptr{b} -- (c);
    \draw[seta] \ptr{c} -- (n);
    \draw[setar] \ptr{a} to[out=-60, in=-120] node[below, leg, text=UFSRed] {1} (c.south);
    \node[leg, anchor=west] at (-0.4,-1.55) {\textbf{1.} \texttt{ant->prox = p->prox;}};
    \node[leg, anchor=west] at (-0.4,-1.95) {\textbf{2.} \texttt{free(p);}};
  \end{tikzpicture}
    \end{column}
    \begin{column}{0.4\textwidth}
      \small
      Para tirar o 20, o 10 tem de passar a apontar para o 30. Por isso a
      busca anda com \textbf{dois} ponteiros: \texttt{p}, no nó procurado,
      e \texttt{ant}, um passo atrás.

      \medskip
      Na ordem contrária --- \texttt{free} antes --- \texttt{p->prox} seria
      lido de um bloco já liberado.
    \end{column}
  \end{columns}
  \fonte{Celes et al.\ (2004), cap.~10~\cite{celes2004introducao};
    Cormen et al.\ (2009), §10.2, \textsc{List-Delete}~\cite{cormen2009introduction}}
\end{frame}

\begin{frame}[fragile]{Um \texttt{remove\_valor} mais curto}
  \begin{columns}[c]
    \begin{column}{0.52\textwidth}
\begin{lstlisting}[style=tinyC]
/* remove_errada.c */
No *remove_valor(No *ini, int x)
{
    No *ant = NULL, *p = ini;
    while (p != NULL && p->valor != x) {
        ant = p;
        p = p->prox;
    }
    if (p == NULL)
        return ini;
    ant->prox = p->prox;
    free(p);
    return ini;
}
\end{lstlisting}
    \end{column}
    \begin{column}{0.44\textwidth}
      \begin{block}{Pense antes de passar o slide}
        \small Sumiu o \texttt{if (ant == NULL)}.
        \begin{enumerate}
          \item Remover o 40 (último) funciona?
          \item E remover o 99, que não está lá?
          \item E remover o 10? Quanto vale \texttt{ant} nessa hora?
        \end{enumerate}
      \end{block}
    \end{column}
  \end{columns}
\end{frame}

\begin{frame}[fragile]{Execute e compare}
\begin{lstlisting}[style=shellStyle]
$ ./remove_errada
10 -> 20 -> 30 -> 40 -> NULL
busca 20: achou
busca 25: nao achou
Segmentation fault: 11
$ gcc -g -fsanitize=address remove_errada.c -o remove_errada-asan
$ ./remove_errada-asan
==36265==ERROR: AddressSanitizer: SEGV on unknown address 0x000000000008
==36265==The signal is caused by a WRITE memory access.
==36265==Hint: address points to the zero page.
    #0 0x0001041f8e3c in remove_valor remove_errada.c:50
    #1 0x0001041f925c in main remove_errada.c:80
\end{lstlisting}
  \begin{center}
    \scriptsize\textcolor{UFSGray}{A linha 80 é \texttt{lista =
    remove\_valor(lista, 10);} --- a primeira remoção.}
  \end{center}
\end{frame}

\begin{frame}{A resposta: o primeiro nó não tem anterior}
  \begin{columns}[c]
    \begin{column}{0.48\textwidth}
      \begin{center}
      \begin{tikzpicture}[node distance=0.8cm]
        \node[var] (l) {lista};
        \node[nolivre, right=0.9cm of l] (a) {10 \nodepart{two} \phantom{xx}};
        \node[no, right=of a] (b) {20 \nodepart{two} \phantom{xx}};
        \node[nulo, right=0.45cm of b] (n) {\ldots};
        \node[varr, above=0.5cm of a] (p) {p};
        \node[varr, minimum width=1.6cm, above=0.5cm of l] (ant) {ant = NULL};
        \draw[setar] (p) -- (a);
        \draw[setag, dashed] (l) -- (a);
        \draw[seta] \ptr{a} -- (b);
        \draw[seta] \ptr{b} -- (n);
        \draw[setar] (l.south) to[out=-35, in=-145] (b.south);
        \node[leg, text=UFSRed, align=center] at ($(a.south)+(0,-1.25)$)
          {a seta que precisa mudar\\ é a de \texttt{lista}};
      \end{tikzpicture}
      \end{center}
    \end{column}
    \begin{column}{0.48\textwidth}
      \small
      Se o 10 é achado de cara, o laço não roda e \texttt{ant} continua
      \texttt{NULL}. \texttt{ant->prox} escreve em
      \texttt{NULL} $+ 8$ --- o endereço \texttt{0x8} do relatório, pois
      \texttt{prox} fica 8 bytes depois do início do nó.

      \medskip
      Quem aponta para o primeiro nó não é um \texttt{prox}: é a variável
      \texttt{lista} do \texttt{main}. Por isso a função muda \texttt{ini}
      e \textbf{devolve} o novo início.
    \end{column}
  \end{columns}
  \fonte{King (2008), §17.5~\cite{king2008c};
    Serebryany et al.\ (2012)~\cite{serebryany2012asan};
    medido com \texttt{remove\_errada.c}}
\end{frame}

\begin{frame}{Recapitulando: remover}
  \begin{block}{O que aprendemos}
    \begin{itemize}
      \item remover é desviar a seta do anterior e só então liberar o nó;
      \item a busca anda com dois ponteiros: \texttt{ant} um passo atrás de
        \texttt{p};
      \item três casos: não achou, achou o primeiro, achou outro;
      \item remover o primeiro muda o início --- quem chama precisa guardar
        o ponteiro devolvido.
    \end{itemize}
  \end{block}
\end{frame}

% ==================================================================
\section{Ponteiro para ponteiro}
% ==================================================================

\begin{frame}{O caso especial é sempre o mesmo}
  \begin{columns}[T]
    \begin{column}{0.5\textwidth}
      \small
      Na inserção e na remoção, o \texttt{if} extra existe pela mesma razão:
      a seta que precisa mudar às vezes é a variável \texttt{lista},
      às vezes o \texttt{prox} de um nó.

      \medskip
      Mas as duas coisas têm o \textbf{mesmo tipo}: \texttt{No *}. Um
      ponteiro do tipo \texttt{No **} pode apontar para qualquer uma delas.
    \end{column}
    \begin{column}{0.46\textwidth}
      \begin{block}{Ponteiro para ponteiro}
        \small Uma variável \texttt{No **pp} guarda o endereço de um
        \texttt{No *}. \texttt{*pp} é esse ponteiro --- dá para ler e
        \emph{mudar} para onde ele aponta.
      \end{block}
      \small Em vez de ``o nó anterior'', a busca guarda ``o endereço da
      seta que chega ao nó atual''.
    \end{column}
  \end{columns}
  \fonte{King (2008), §17.6~\cite{king2008c};
    Kernighan \& Ritchie (1989), §5.6~\cite{kernighan1989c}}
\end{frame}

\begin{frame}[fragile]{\texttt{ponteiro\_duplo.c}: sem \texttt{if} para o primeiro}
  \begin{columns}[T]
    \begin{column}{0.5\textwidth}
\begin{lstlisting}[style=codC]
void insere_ordenado(No **pp, int x)
{
    while (*pp != NULL && (*pp)->valor < x)
        pp = &(*pp)->prox;
    No *novo = malloc(sizeof *novo);
    if (novo == NULL) {
        perror("malloc");
        exit(1);
    }
    novo->valor = x;
    novo->prox = *pp;
    *pp = novo;
}
\end{lstlisting}
    \end{column}
    \begin{column}{0.46\textwidth}
\begin{lstlisting}[style=codC]
void remove_valor(No **pp, int x)
{
    while (*pp != NULL && (*pp)->valor != x)
        pp = &(*pp)->prox;
    if (*pp != NULL) {
        No *p = *pp;
        *pp = p->prox;
        free(p);
    }
}
\end{lstlisting}
\begin{lstlisting}[style=codC]
/* no main */
insere_ordenado(&lista, 30);
remove_valor(&lista, 10);
\end{lstlisting}
    \end{column}
  \end{columns}
\end{frame}

\begin{frame}{Pense antes de passar o slide}
  \begin{block}{Leia o código com calma}
    \begin{enumerate}
      \item Por que o \texttt{main} passa \texttt{\&lista}, e não
        \texttt{lista}? E por que as funções agora são \texttt{void}?
      \item No começo, para onde \texttt{pp} aponta? E depois de uma volta
        do laço?
      \item Por que \texttt{(*pp)->valor} precisa dos parênteses? E o que
        é \texttt{\&(*pp)->prox}?
      \item Remover o 10 (primeiro) e o 40 (último): siga \texttt{pp} nos
        dois casos.
    \end{enumerate}
  \end{block}
\end{frame}

\begin{frame}{A resposta: \texttt{pp} aponta para a seta, não para o nó}
  \begin{center}
  \begin{tikzpicture}[node distance=1.5cm]
    \node[var] (l) {lista};
    \node[no, right=1.2cm of l] (a) {10 \nodepart{two} \phantom{xx}};
    \node[no, right=of a] (b) {20 \nodepart{two} \phantom{xx}};
    \node[no, right=of b] (c) {30 \nodepart{two} \phantom{xx}};
    \node[nulo, right=0.6cm of c] (n) {NULL};
    \draw[seta] (l) -- (a);
    \draw[seta] \ptr{a} -- (b);
    \draw[seta] \ptr{b} -- (c);
    \draw[seta] \ptr{c} -- (n);
    \node[varr] (p0) at ($(l)+(0,-1.0)$) {pp};
    \draw[setar] (p0) -- (l);
    \node[leg, below=0.05cm of p0, align=center] {início\\ \texttt{\&lista}};
    \node[varr] (p1) at ($(a.two south)+(0,-1.0)$) {pp};
    \draw[setar] (p1) -- (a.two south);
    \node[leg, below=0.05cm of p1, align=center] {1 volta\\ \texttt{\&}(prox do 10)};
    \node[varr] (p2) at ($(b.two south)+(0,-1.0)$) {pp};
    \draw[setar] (p2) -- (b.two south);
    \node[leg, below=0.05cm of p2, align=center] {2 voltas\\ \texttt{\&}(prox do 20)};
  \end{tikzpicture}
  \end{center}
  \small
  \texttt{*pp = p->prox} muda \emph{a seta para onde \texttt{pp} aponta} ---
  seja ela \texttt{lista} ou um \texttt{prox}. Por isso não há caso
  especial, e a função não precisa devolver nada: ela mexe no
  \texttt{lista} do \texttt{main} pelo endereço. Os parênteses são
  necessários porque \texttt{->} tem precedência maior que \texttt{*} e
  \texttt{\&}: \texttt{\&(*pp)->prox} é ``o endereço do campo
  \texttt{prox} do nó atual''.
  \fonte{King (2008), §17.6~\cite{king2008c};
    Kernighan \& Ritchie (1989), §2.12~\cite{kernighan1989c};
    medido com \texttt{ponteiro\_duplo.c}}
\end{frame}

\begin{frame}{Recapitulando: ponteiro para ponteiro}
  \begin{block}{O que aprendemos}
    \begin{itemize}
      \item o caso especial vem de a seta do primeiro nó morar fora da
        lista, na variável \texttt{lista};
      \item \texttt{No **pp} aponta para uma seta qualquer, de dentro ou de
        fora da lista;
      \item avançar é \texttt{pp = \&(*pp)->prox}; mudar a seta é
        \texttt{*pp = \ldots};
      \item quem chama passa \texttt{\&lista}, e a função não devolve o
        início.
    \end{itemize}
  \end{block}
\end{frame}

% ==================================================================
\section{Quanto custa cada operação}
% ==================================================================

\begin{frame}{Custo das operações numa lista de $n$ nós}
  \begin{center}
    \footnotesize
    \begin{tabular}{@{}lll@{}}
      \toprule
      operação & nós visitados (pior caso) & por quê\\
      \midrule
      buscar & $n$ & só se lê na ordem das setas\\
      inserir no início & constante & não percorre\\
      inserir no fim & $n$ (constante guardando o último) & acha o último\\
      inserir em ordem & $n$ & acha o anterior\\
      remover um valor & $n$ & busca, depois desvio\\
      remover um nó já conhecido & $n$ & precisa achar o anterior\\
      \bottomrule
    \end{tabular}
  \end{center}
  \small
  Encaixar e desligar um nó é barato: dois ponteiros. O caro é
  \textbf{chegar} ao lugar. A lista ordenada não ajuda tanto quanto o vetor
  ordenado: sem acesso por índice, não há busca binária.
  \fonte{Cormen et al.\ (2009), §10.2~\cite{cormen2009introduction};
    Knuth (1997), §2.2.3~\cite{knuth1997art}}
\end{frame}

\begin{frame}{Recapitulando: quanto custa cada operação}
  \begin{block}{O que aprendemos}
    \begin{itemize}
      \item ligar e desligar um nó custa o mesmo com 3 ou 3 milhões de nós;
      \item achar a posição custa até $n$ passos;
      \item uma lista ordenada não permite busca binária;
      \item remover um nó conhecido ainda exige achar o anterior --- o
        problema que a próxima aula resolve.
    \end{itemize}
  \end{block}
\end{frame}

% ==================================================================
\section{Mãos à obra}
% ==================================================================

\begin{frame}{Altere o programa}
  \small
  \begin{block}{Tarefa 1}
    A lista está em ordem: ao passar de \texttt{x}, a busca já pode
    desistir. Altere \texttt{busca} para isso. Quantos nós ela visita agora
    em \texttt{busca(lista, 25)}? E antes?
  \end{block}
  \begin{block}{Tarefa 2}
    Insira o 20 uma segunda vez. O novo 20 entra antes ou depois do antigo?
    Troque \texttt{<} por \texttt{<=} no laço e preveja o que muda antes de
    executar.
  \end{block}
  \begin{block}{Tarefa 3}
    Em \texttt{ponteiro\_duplo.c}, escreva \texttt{void remove\_todos(No
    **pp, int x)}, que tira \emph{todas} as ocorrências de \texttt{x}. Teste
    com \texttt{10 -> 10 -> 20 -> 10}.
  \end{block}
\end{frame}

\begin{frame}{Exercícios (1/2)}
  \footnotesize
  \begin{block}{Exercício 1 --- inserir no fim sem percorrer}
    Crie \texttt{typedef struct \{ No *ini, *fim; \} Lista;} e escreva
    \texttt{insere\_fim(Lista *l, int x)} com custo constante. Que funções
    precisam atualizar \texttt{fim}? Refaça a tabela de tempos.
  \end{block}
  \begin{block}{Exercício 2 --- remover por posição}
    \texttt{No *remove\_pos(No *ini, int k)} remove o nó da posição
    \texttt{k} (a partir de 0), se existir. Escreva também a versão com
    \texttt{No **}.
  \end{block}
  \begin{block}{Exercício 3 --- busca recursiva}
    \texttt{No *busca\_rec(No *ini, int x)}. Quais são os dois casos base?
    Com 1 milhão de nós, ela funciona sem \texttt{-O2}? E com?
  \end{block}
  \fonte{Feofiloff (2009), cap.~4~\cite{feofiloff2009algoritmos};
    Celes et al.\ (2004), cap.~10~\cite{celes2004introducao}}
\end{frame}

\begin{frame}{Exercícios (2/2)}
  \footnotesize
  \begin{block}{Exercício 4 --- mover para a frente}
    \texttt{No *busca\_mtf(No *ini, int x)}: ao achar \texttt{x}, move o
    nó para o início e devolve o novo início. Quais ponteiros mudam?
  \end{block}
  \begin{block}{Exercício 5 --- intercalar}
    \texttt{No *intercala(No *a, No *b)} junta duas listas ordenadas numa
    só, ordenada, \textbf{sem} \texttt{malloc}: só religando os nós
    existentes.
  \end{block}
  \begin{block}{Exercício 6 --- inverter no lugar}
    \texttt{No *inverte(No *ini)} inverte a lista mudando apenas os
    \texttt{prox}, sem criar nós. Quantos ponteiros auxiliares são
    necessários?
  \end{block}
  \fonte{King (2008), §17.5--17.6~\cite{king2008c};
    Tenenbaum et al.\ (1995), cap.~4~\cite{tenenbaum1995estruturas}}
\end{frame}

\begin{frame}{Desafio}
  \begin{block}{Uma agenda sempre em ordem alfabética}
    \small Leia comandos da entrada padrão até o fim do arquivo e mantenha
    uma lista de
    \begin{center}
      \texttt{typedef struct contato \{ char *nome; struct contato *prox; \} Contato;}
    \end{center}
    \vspace{-0.6em}
    \begin{itemize}
      \item \texttt{+ nome} insere em ordem alfabética (\texttt{strcmp}),
        sem repetir nomes;
      \item \texttt{- nome} remove; \texttt{? nome} responde ``sim'' ou
        ``não''; \texttt{=} lista todos;
      \item cada nome mora num bloco do \texttt{malloc} com o tamanho exato
        --- e é liberado junto com o nó.
    \end{itemize}
    Use \texttt{Contato **} para não ter caso especial. Com
    \texttt{-fsanitize=address}, nenhum erro e nenhum vazamento.
  \end{block}
  \fonte{King (2008), §17.5 e §17.6~\cite{king2008c}}
\end{frame}

\begin{frame}{Fechamento}
  \begin{columns}[T]
    \begin{column}{0.5\textwidth}
      \textbf{O que mudou hoje}\\[0.3em]
      \small A lista deixou de crescer só pela frente: agora dá para achar
      um valor, encaixar um nó em qualquer ponto e tirar um nó do meio.
    \end{column}
    \begin{column}{0.5\textwidth}
      \textbf{O que ficou}\\[0.3em]
      \small Para mexer num nó, é preciso o \emph{anterior} --- e a lista só
      anda para a frente.
    \end{column}
  \end{columns}
  \begin{block}{Para levar}
    \begin{itemize}
      \small
      \item buscar: \texttt{p != NULL \&\& p->valor != x}, nessa ordem;
      \item inserir: \texttt{novo->prox = ant->prox; ant->prox = novo;}
      \item remover: \texttt{ant->prox = p->prox; free(p);}
      \item o primeiro nó é especial --- a menos que se use \texttt{No **}.
    \end{itemize}
  \end{block}
  \small\textcolor{UFSGray}{Próxima aula: um nó que aponta também para o
  anterior --- a lista duplamente encadeada.}
\end{frame}

\section*{Referências}
\begin{frame}[allowframebreaks]{Referências}
  \nocite{*}
  \bibliographystyle{plain}
  \bibliography{referencias}
\end{frame}

\end{document}
